\section{A fifth-iterate obstruction for the theta kernel}
Set $\mathcal Tf=(f')^2-ff''$ and
\begin{align*}
\Phi(r)&=\sum_{m\ge1}q_m(2q_me^{4r}-3)e^{5r-q_me^{4r}},\\
q_m&=\pi m^2,\qquad s(t)=\Phi(\sqrt t).
\end{align*}
The theta-functional equation makes $\Phi$ even; thus $s$ is analytic
near zero. This is the kernel normalization used by
Planat--Sol\'e~\cite{PS}; constant and argument rescalings do not
change the signs under consideration.

\begin{theorem}[Certified obstruction]\label{thm:five}
One has
\[
-4.813\cdot 10^{65}<\mathcal T^5s(0)
<-4.812\cdot10^{65}
\]
and $\mathcal T^js(0)>0$ for $0\le j\le4$. In particular
$\mathcal T^4s$ is positive but strictly log-convex at zero,
and $\mathcal T^5s<0$ on some punctured interval to the right.
\end{theorem}
The second-level concavity theorem of Planat--Sol\'e~\cite{PS}
is not contradicted by this higher-order local obstruction.

\begin{lemma}[Finite theta tail]
For each $0\le k\le10$, truncation at $m=12$ in the exact
Taylor-coefficient formula
\begin{align*}
a_k&=[t^k]s(t)\\
&=\frac1{(2k)!}\sum_{m\ge1}q_me^{-q_m}P_{2k}(q_m),\\
P_0(x)&=2x-3,\\
P_{j+1}(x)&=(5-4x)P_j(x)+4xP_j'(x),
\end{align*}
has absolute error below $4\cdot10^{-83}$.
\end{lemma}
\begin{proof}
Put $C_j=\sum|\operatorname{coeff}(P_j)|$.
For $x\ge1$, $|P_{2k}(x)|\le C_{2k}x^{2k+1}$.
Using $3<\pi<22/7$ and $e>2$, the omitted $m$-th
summand has magnitude at most
\[
\frac{C_{2k}}{(2k)!}
\left(\frac{22m^2}7\right)^{2k+2}2^{-3m^2}.
\]
The successive ratio for $m\ge13$ and $0\le k\le10$ is
bounded by $(14/13)^{44}2^{-81}<1/2$.
Thus the absolute tail is at most
\[
B_k=\frac{2C_{2k}}{(2k)!}
\left(\frac{22\cdot169}7\right)^{2k+2}2^{-507}.
\]
Exact integer comparisons in the supplied Python certificate verify
$B_k<4\cdot10^{-83}$ for all eleven values of $k$.
\end{proof}

\begin{proof}[Proof of Theorem~\ref{thm:five}]
Only $a_0,\ldots,a_{10}$ are required. If $f=\sum b_kt^k$ then
\begin{align*}
[t^k]\mathcal Tf
&=\sum_{i+j=k}(i+1)(j+1)b_{i+1}b_{j+1}\\
&\quad-\sum_{i+j=k}(j+1)(j+2)b_ib_{j+2}.
\end{align*}
Five successive applications consume exactly ten trailing Taylor
coefficients. To certify these initial coefficients, enclose
$\pi=16\arctan(1/5)-4\arctan(1/239)$ using alternating rational
series with 100 and 40 terms. Enclose $e^{-q_m}$ for $1\le m\le12$
via the reciprocal of $[e^{q_m/1024}]^{1024}$: Taylor degree 80
for the reduced argument has a rigorous omitted bound
$2\cdot2^{-81}/81!<10^{-110}$.
Every arithmetic operation is an outward-rounded integer
operation on intervals with denominator $10^{110}$. Expand
the input coefficient intervals by the analytic theta-tail bound
and propagate them five times by the displayed recurrence.
The source \texttt{certificates/xi\_fifth\_iterate\_exact.py}
reproduces the stated enclosure without third-party packages
or floating-point decisions. The signs for levels $0$ to $4$ are
strictly positive, and level $5$ strictly negative.
Finally $(\log\mathcal T^4s)''(0)
=-\mathcal T^5s(0)/(\mathcal T^4s(0))^2>0$, and analyticity
gives the asserted neighborhood.
\end{proof}
